% ---- pg-478 （书页 466） ----

表观失误率为
\[
\mathrm{APER}=\left(\frac{2+2}{12+12}\right)\times100\%=\left(\frac{4}{24}\right)\times100\%=16.7\%
\]
\hfill$\blacksquare$

APER 具有直观上的感染力, 且易于计算. 不幸的是, 它倾向于低估 AER, 而且除非 $n_1$ 和 $n_2$ 非常大, 这个问题不会消失. 出现这种乐观估计的原因, 本质上由于是被用来构造分类函数的数据又被用于对这个函数进行评估.

可以构造比表观失误率更好的失误率估计值, 而且这种估计值相对来说也比较容易计算, 且不需要对分布做假定. 一种方法是将整个样本分成两部分: 训练样本和验证样本. 训练样本用于构造分类函数, 而验证样本则用于评估. 失误率用验证样本中被错分数据的比例来确定. 虽然这种方法不用同一组数据既构造又评估分类函数, 从而克服了偏倚问题, 但它有两个主要缺陷:

(i) 需要用大样本.

(ii) 所评估的函数并非感兴趣的函数. 这是因为最终几乎所有的数据都必须用于构造分类函数. 如果不是这样, 就会损失有价值的信息.

还有一种方法看起来效果较好, 称为拉亨布鲁克(Lachenbruch)的“提留”方法\footnote{此方法有时称为“刀切”, 或“交叉验证”.} (见文献[24]):

1. 从 $\pi_1$ 组观测值开始, 从这组数据中取出一个, 用其余 $n_1-1$ 个观测值和 $\pi_2$ 组的 $n_2$ 个观测值构造分类函数.

2. 用第 1 步中所构造的函数对“提留”的观测值进行分类.

3. 重复步骤 1 和 2, 直至 $\pi_1$ 中所有 $n_1$ 个观测值都被分类完毕为止. 令 $n_{1M}^{(H)}$ 为 $\pi_1$ 组中被错分的提留观测值个数.

4. 对 $\pi_2$ 中的观测值重复步骤 1 至 3. 令 $n_{2M}^{(H)}$ 为 $\pi_2$ 组中被错分的提留观测值个数.

在\hyperref[eq:11-1]{式(11-1)}和\hyperref[eq:11-2]{式(11-2)}中的条件错分概率 $P(2\mid1)$ 和 $P(1\mid2)$ 的估计量为:
\phantomsection\label{eq:11-35}
\begin{equation}
\begin{aligned}
\hat{P}(2\mid1)&=\frac{n_{1M}^{(H)}}{n_1}\\
\hat{P}(1\mid2)&=\frac{n_{2M}^{(H)}}{n_2}
\end{aligned}
\tag{11-35}
\end{equation}
总错分比例 $(n_{1M}^{(H)}+n_{2M}^{(H)})/(n_1+n_2)$ 为期望真实失误率 $E(\mathrm{AER})$ 的一个接近无偏的估计量, 即使对中等样本这个结论也对.

\begin{mvbox}
\phantomsection\label{eq:11-36}
\begin{equation}
\hat{E}(\mathrm{AER})=\frac{n_{1M}^{(H)}+n_{2M}^{(H)}}{n_1+n_2}
\tag{11-36}
\end{equation}
\end{mvbox}

拉亨布鲁克的提留方法与线性分类函数共同使用时, 计算上也是可行的. 在某些现有的判别分析的计算机程序中, 这种方法是所提供选择项目之一.

\noindent\textbf{例 11.7}\quad (用提留方法计算失误率的估计值)

我们将对\hyperref[eq:11-18]{式(11-18)}的相同代价和相同先验概率这种情况来说明拉亨布鲁克的提留方法和失误率估计值的计算. 考虑下列数据矩阵和描述性统计数字. (我们假定 $n_1=n_2=3$ 个二元观测值随机地取自两个总体 $\pi_1$ 和 $\pi_2$, 并假定它们有公共协方差阵.)
